%-------------------------------------------------------------------------------
%  PERSONAL INFORMATION
%  Comment any of the lines below if they are not required
%-------------------------------------------------------------------------------
% Available options: circle|rectangle,edge/noedge,left/right
% \photo{profile}
\name{Arnaud}{Tanguy}
\position{Ingénieur de recherche robotique{\enskip\cdotp\enskip}Expert C++}
\address{Résidence les Béguines, Bât B., 3035 Avenue des Moulins, 34080 Montpellier, France}

\mobile{0762705579}
\email{arn.tanguy@gmail.com}
\homepage{arntanguy.fr}
\github{arntanguy}
\linkedin{arnaud-tanguy}
% \gitlab{gitlab-id}
% \stackoverflow{SO-id}{SO-name}
% \twitter{@twit}
% \skype{skype-id}
% \reddit{reddit-id}
% \medium{madium-id}
\googlescholar{SeWUqycAAAAJ}{Google Scholar}
%% \firstname and \lastname will be used
% \googlescholar{googlescholar-id}{}
% \extrainfo{extra informations}

% \quote{Ingénieur de recherche en robotique logicielle, expert en développement C++ et vision par ordinateur
% avec 8 ans d'expérience (LIRMM, JRL, I3S). Co-développeur et mainteneur du framework open-source mc\_rtc,
% je suis spécialisé dans le contrôle de systèmes robotiques complexes (humanoïdes, quadrupèdes) et
% l'intégration de pipelines DevOps/build rigoureux pour sécuriser le déploiement sur cibles réelles. Engagé
% pour la reproductibilité et la science ouverte.}
% \quote{Ingénieur de recherche en robotique, expert en C++, fort de 8 ans d'expérience en robotique
% humanoïde et vision (LIRMM, JRL). Co-développeur de mc\_rtc et fervent défenseur de la science ouverte, je
% m’anime à concevoir des architectures logicielles et pipelines DevOps robustes pour propulser la recherche
% du laboratoire sur des robots réels.}
% \quote{Ingénieur de recherche en robotique logicielle, expert en développement C++ et vision par ordinateur
% avec 8 ans d'expérience (LIRMM, JRL, I3S). Ma force ? Allier l'exigence de la science ouverte
% (co-développeur de mc\_rtc) à la puissance du DevOps pour transformer des concepts complexes en systèmes
% robustes, directement déployables sur robots réels et immédiatement exploitables pour les travaux des chercheurs.}
% \quote{Fort de mon expertise en développement C++, robotique humanoïde et vision par ordinateur, avec 12
% ans d'expérience en laboratoire (LIRMM, JRL, I3S), je m'attache à mettre mes compétences à l'appui des
% travaux des laboratoires. Je suis particulièrement attaché aux concepts de science ouverte, de
% reproductibilité des travaux, et m'attache à mettre en place les outils permettant à la fois de péréniser
% les travaux, mais aussi de rendre les recherches en cours possible.}

% \quote{Fort de 12 ans d'expérience en laboratoire (LIRMM, JRL, I3S), je mets mon expertise en développement
% C++, robotique humanoïde et vision par ordinateur au service direct des travaux de recherche. Fervent
% défenseur de la science ouverte et de la reproductibilité, je m’investis pleinement dans le déploiement
% d’outils et d’architectures DevOps robustes. Mon objectif est double : pérenniser les contributions passées
% et fournir le cadre permettant le développement efficace des recherches futures.}

% \quote{Fort de 12 ans d'expérience en laboratoire (LIRMM, JRL, I3S), je mets mon expertise en développement
% C++, robotique humanoïde et vision par ordinateur au service direct des travaux de recherche. Fervent
% défenseur de la science ouverte et de la reproductibilité, je m’investis pleinement dans le déploiement
% d'architectures logicielles fiables et d'environnements de développement robustes. Mon objectif est double
% : pérenniser les contributions passées et fournir le cadre permettant le développement efficace des
% recherches futures.}

% \quote{Ingénieur de Recherche spécialisé dans la mise en œuvre et la conduite expérimentale sur robots
% réels (humanoïdes, bras manipulateurs, capteurs tiers). Fort de 12 ans de pratique terrain au sein de
% laboratoires internationaux, j'articule mon expertise en développement C++ autour des défis de
% l'instrumentation et du maintien en conditions opérationnelles des plateformes. Mon objectif est de faire
% le pont entre la théorie et le matériel, en garantissant des environnements de test robustes, sécurisés et
% reproductibles pour la recherche.}

% 1. Define a clean macro for a subtle, charcoal bold text style
\newcommand{\keyword}[1]{{\textsc{#1}}}
